\section{Phase 6 推导草稿：Qualified Contract-Manufacturing Market 与 Labor-Market Clearing}

本阶段开启 benchmark MAH 模型中的 market-clearing block，但这里的任务严格限定在 setup 层面，而不是完整的 stationary competitive equilibrium 证明。此时我们把 Phase 2 已经完成的优化后静态控制，特别是 $x_B^*(z;p_m,M)$ 与 $m^*(z;p_m,w)$；Phase 3 已经完成的 fixed-price value-function block；Phase 5 已经完成的 entrant masses $n_E^A(M),n_E^B(M),n_E^C(M)$；紧状态空间 $Z$；以及此前已经声明过的正则条件，都视为给定输入。我们还把 $\mu_A,\mu_B,\mu_C$ 视为给定的 candidate finite Borel measures，也就是 candidate incumbent measures。在 Phase 6 的 opening setup 中，它们还不能被当作 stationary objects，因为 invariant-measure 论证属于后续 Phase 7，而不属于这里。

\subsection{Benchmark Fidelity Check}

\paragraph{Step 1：Phase 6 相关的精确 benchmark 对象。}
Phase 6 opening setup 中真正相关的对象，全部来自 benchmark manuscript 与此前已经完成的 Phase 1--5 逻辑。来自 Phase 2 的对象包括优化后控制 $x_B^*(z;p_m,M)$ 与 $m^*(z;p_m,w)$，以及已经推导完成、但这里不会被重新压缩写成 generic $\Pi_j$ 的优化后 payoff objects $\Pi_A(z;w,M)$、$\Pi_B(z;p_m,M)$、$\Pi_C(z;p_m,w,M)$。来自 Phase 3 的对象是已经完成的 fixed-price value-function block，但这里仅把它当作后续阶段已经可用的动态背景，而不在本节重做推导。来自 Phase 5 的对象是 entrant masses $n_E^A(M),n_E^B(M),n_E^C(M)$。除此之外，本阶段还把紧状态空间 $Z$、candidate incumbent measures $\mu_A,\mu_B,\mu_C$、reduced-form commercialization-success object $\lambda_B(z,p_m,M)$、劳动对象 $\ell_A(z;w)$、$\ell_B(z;p_m,M)$、$\ell_C(z;p_m,w)$、进入者启动劳动系数 $\ell_E^j$ 以及外生劳动供给 $\bar L$ 视为相关 benchmark 对象。

\paragraph{Step 2：精确的 benchmark 市场对象与 normalization。}
contract-manufacturing block 使用一个必须在这里明说的 benchmark normalization：one successfully commercialized type-$B$ original-drug idea requires one unit of qualified external manufacturing service。这个表述只是 benchmark 市场方程中的 normalization，不是新的 implementation block，不是新的 contracting layer，也不是新的 primitive optimization problem。在这个 normalization 下，合格受托生产市场的需求与供给分别定义为
\begin{align}
D_m(p_m,M)
&=\int_Z x_B^*(z;p_m,M)\,\lambda_B(z,p_m,M)\,\mu_B(dz),
\label{eq:p6zh_Dm}\\
S_m(p_m,w)
&=\int_Z m^*(z;p_m,w)\,\mu_C(dz),
\label{eq:p6zh_Sm}
\end{align}
市场出清条件是
\begin{equation}
D_m(p_m,M)=S_m(p_m,w).
\label{eq:p6zh_pm_clear}
\end{equation}

benchmark labor-demand object 直接写为
\begin{equation}
\begin{aligned}
L_D(M)={}&\int_Z \ell_A(z;w)\,\mu_A(dz)
+\int_Z \ell_B(z;p_m,M)\,\mu_B(dz)\\
&+\int_Z \ell_C(z;p_m,w)\,\mu_C(dz)
+\sum_{j\in\{A,B,C\}}\ell_E^j n_E^j(M),
\end{aligned}
\end{equation}
benchmark 的劳动市场出清条件是
\begin{equation}
L_D(M)=\bar L.
\end{equation}

\paragraph{Step 3：本阶段此时能够忠实完成什么，以及还不能声称什么。}
在 opening setup 这一步，Phase 6 能够忠实完成的工作，是在给定 candidate prices、candidate incumbent measures 与 entrant masses 的条件下，把 benchmark 的两个市场对象精确重述出来，并说明这些市场对象调用的是哪些前序 phase 已经完成的 policy objects。但是，这一步还不能声称 candidate incumbent measures 已经是 invariant measures，还不能声称 full stationary competitive equilibrium existence 已经得到证明，也还不能声称 prices 与 measures 的 joint fixed point 已经被求出。这些内容都超出了 Phase 1--5 已完成逻辑在当前时点能够支持的范围，因此必须明确标记为 pending，而不能在这里静默假定。

\paragraph{Step 4：为什么必须先写 Phase 6，再写 Phase 7 与 Phase 8。}
必须先把 Phase 6 单独写出来，因为后续 stationary-distribution block 与 full-equilibrium block 都需要一个清楚、可核对的市场接口。Phase 7 若要讨论 stationary consistency，就必须先知道 candidate incumbent measures 最终要和哪两个 benchmark market-clearing equations 同时成立。Phase 8 若要讨论 full stationary competitive equilibrium，就必须把 value-function block、entry block、market-clearing block 与 stationary-distribution block 分开写清楚后再组合起来。所以，Phase 6 的角色是先隔离出 conditional market-clearing block，而不是提前把 invariant measures 或 joint price-measure closure 一并宣称完成。

因此，Phase 6 的后续部分将分别在给定 candidate prices、candidate incumbent measures 与 entrant masses 的条件下，推导 contract-manufacturing market object 与 labor-market object。那些后续推导仍然只是 conditional market-block calculations，而不是 invariant measures、joint price-measure closure 或 full stationary competitive equilibrium existence 的证明。

\subsection{Part 1：合格受托生产市场}

\paragraph{Benchmark 对象重述。}
在 benchmark manuscript 中，合格受托生产市场由 aggregate demand object
\[
D_m(p_m,M)=\int_Z x_B^*(z;p_m,M)\,\lambda_B(z,p_m,M)\,\mu_B(dz),
\]
aggregate supply object
\[
S_m(p_m,w)=\int_Z m^*(z;p_m,w)\,\mu_C(dz),
\]
以及 market-clearing restriction
\[
D_m(p_m,M)=S_m(p_m,w)
\]
来刻画。这些都是 benchmark 的原始市场对象，本节保持原样，不做 redesign。本阶段中的 $\mu_B$ 与 $\mu_C$ 只是 candidate finite Borel measures，还没有被证明为 invariant stationary measures。控制变量 $x_B^*(z;p_m,M)$ 与 $m^*(z;p_m,w)$ 都是从 Phase 2 继承来的、已经求解完成的 policy objects，本节不重新打开它们的静态优化问题。

\paragraph{经济作用。}
这里必须把 benchmark normalization 明确说出来：one successfully commercialized type-$B$ original-drug idea requires one unit of qualified external manufacturing service。正是这个 normalization 解释了为什么需求对象中直接出现乘积 $x_B^*(z;p_m,M)\lambda_B(z,p_m,M)$。对于状态为 $z$ 的 type-$B$ firm，$x_B^*(z;p_m,M)$ 是最优 original-drug idea-generation intensity，$\lambda_B(z,p_m,M)$ 是这些 ideas 通过 qualified external manufacturing 成功商业化的 reduced-form success rate。因此，它们的乘积表示该 firm-state 对 qualified external manufacturing services 的期望需求量。再对 candidate type-$B$ measure $\mu_B$ 积分，就得到由 type-$B$ firms 产生的 expected industry demand for qualified external manufacturing。

供给对象在 type-$C$ 一侧具有平行的解释。对于状态为 $z$ 的 type-$C$ firm，$m^*(z;p_m,w)$ 是在候选价格 $(p_m,w)$ 下最优选择的 qualified contract-manufacturing supply。对 candidate type-$C$ measure $\mu_C$ 积分，就得到由 type-$C$ incumbents 提供的 industry supply of qualified external manufacturing。因此，这个市场比较的是两边的 aggregate objects：一边是 type-$B$ commercialization 所引致的 outsourcing demand，另一边是 type-$C$ incumbents 所提供的 qualified manufacturing supply。

\paragraph{数学问题。}
固定制度 $M\in\{0,1\}$、候选的 qualified-manufacturing price $p_m$、候选工资 $w$，以及 candidate finite Borel measures $\mu_B$ 与 $\mu_C$。本部分的数学任务只限于三个问题。第一，定义 $D_m(p_m,M)$ 与 $S_m(p_m,w)$ 的两个积分是否良定义且有限。第二，在价格的局部紧域上附加连续性条件后，这两个 aggregate objects 是否会随 prices 连续变化。第三，在这两个对象已经构造出来之后，如何把 benchmark market-clearing restriction 记录成一个 excess-demand object，供后续均衡分析使用。本部分不求解该 restriction，不声称 clearing price 已经存在，也不声称 clearing price 唯一。

\paragraph{Supplemental assumptions.}
\begin{assumption}[可测性与候选测度有限性]\label{ass:p6:cm-meas}
对每个固定的三元组 $(p_m,w,M)$，映射
\[
z\mapsto x_B^*(z;p_m,M)\lambda_B(z,p_m,M)
\]
在 $Z$ 上 Borel 可测，映射
\[
z\mapsto m^*(z;p_m,w)
\]
在 $Z$ 上 Borel 可测。候选测度 $\mu_B$ 与 $\mu_C$ 都是定义在 $Z$ 上的有限 Borel measures。
\end{assumption}

\begin{assumption}[固定候选价格下的有界性]\label{ass:p6:cm-bound}
对每个固定的三元组 $(p_m,w,M)$，存在有限常数 $B_D(p_m,M)$ 与 $B_S(p_m,w)$，使得
\[
\left|x_B^*(z;p_m,M)\lambda_B(z,p_m,M)\right|\le B_D(p_m,M)
\quad\text{对所有 } z\in Z,
\]
并且
\[
\left|m^*(z;p_m,w)\right|\le B_S(p_m,w)
\quad\text{对所有 } z\in Z.
\]
\end{assumption}

\begin{assumption}[紧局部价格域上的联合连续性]\label{ass:p6:cm-cont}
固定制度 $M$。令 $\mathcal P_m^{\mathrm{loc}}\subset\mathbb R_+$ 表示 qualified-manufacturing price 的一个紧局部价格域，令 $\mathcal W^{\mathrm{loc}}\subset\mathbb R_+$ 表示工资的一个紧局部价格域。映射
\[
(z,p_m)\mapsto x_B^*(z;p_m,M)\lambda_B(z,p_m,M)
\]
在 $Z\times\mathcal P_m^{\mathrm{loc}}$ 上联合连续，映射
\[
(z,p_m,w)\mapsto m^*(z;p_m,w)
\]
在 $Z\times\mathcal P_m^{\mathrm{loc}}\times\mathcal W^{\mathrm{loc}}$ 上联合连续。
\end{assumption}

\begin{proposition}[良定义性与有限性]
固定制度 $M$、候选价格 $(p_m,w)$ 以及 candidate finite Borel measures $(\mu_B,\mu_C)$。在 Assumption~\\ref{ass:p6:cm-meas} 与 Assumption~\\ref{ass:p6:cm-bound} 下，$D_m(p_m,M)$ 与 $S_m(p_m,w)$ 都良定义，并且都是有限实数。
\end{proposition}

\begin{proof}
固定 $(p_m,w,M,\mu_B,\mu_C)$。定义两个 integrands：
\[
f_D(z;p_m,M):=x_B^*(z;p_m,M)\lambda_B(z,p_m,M)
\]
以及
\[
f_S(z;p_m,w):=m^*(z;p_m,w).
\]
由 Assumption~\\ref{ass:p6:cm-meas}，$f_D(\cdot;p_m,M)$ 与 $f_S(\cdot;p_m,w)$ 都在 $Z$ 上 Borel 可测。因此，它们都可以分别对相应的 candidate Borel measure 进行积分。

接下来使用 Assumption~\\ref{ass:p6:cm-bound}。对任意 $z\in Z$，
\[
|f_D(z;p_m,M)|=\left|x_B^*(z;p_m,M)\lambda_B(z,p_m,M)\right|\le B_D(p_m,M).
\]
把这个不等式对 $\mu_B$ 积分，得到
\[
\int_Z |f_D(z;p_m,M)|\,\mu_B(dz)
\le
\int_Z B_D(p_m,M)\,\mu_B(dz).
\]
由于 $B_D(p_m,M)$ 对 $z$ 来说是常数，右边可化为
\[
\int_Z B_D(p_m,M)\,\mu_B(dz)
=
B_D(p_m,M)\mu_B(Z).
\]
Assumption~\\ref{ass:p6:cm-meas} 说明 $\mu_B$ 是有限 Borel measure，所以 $\mu_B(Z)<\infty$。Assumption~\\ref{ass:p6:cm-bound} 说明 $B_D(p_m,M)<\infty$。因此
\[
\int_Z |f_D(z;p_m,M)|\,\mu_B(dz)<\infty.
\]
这说明 $f_D(\cdot;p_m,M)$ 关于 $\mu_B$ 绝对可积，所以定义 $D_m(p_m,M)$ 的积分存在且有限。

对供给对象作完全平行的论证。对任意 $z\in Z$，
\[
|f_S(z;p_m,w)|=|m^*(z;p_m,w)|\le B_S(p_m,w).
\]
对 $\mu_C$ 积分得到
\[
\int_Z |f_S(z;p_m,w)|\,\mu_C(dz)
\le
\int_Z B_S(p_m,w)\,\mu_C(dz)
=
B_S(p_m,w)\mu_C(Z).
\]
由 Assumption~\\ref{ass:p6:cm-meas}，$\mu_C(Z)<\infty$；由 Assumption~\\ref{ass:p6:cm-bound}，$B_S(p_m,w)<\infty$。因此
\[
\int_Z |f_S(z;p_m,w)|\,\mu_C(dz)<\infty.
\]
这说明定义 $S_m(p_m,w)$ 的积分也存在且有限。

由此我们逐步证明了：在固定 candidate prices 与固定 candidate measures 下，两个 benchmark market objects 都良定义且有限。
\end{proof}

\begin{proposition}[紧局部价格域上的连续性]
固定制度 $M$ 与 candidate finite Borel measures $(\mu_B,\mu_C)$。在 Assumption~\\ref{ass:p6:cm-meas}--Assumption~\\ref{ass:p6:cm-cont} 下，需求对象 $D_m(\cdot,M)$ 在 $\mathcal P_m^{\mathrm{loc}}$ 上连续，供给对象 $S_m(\cdot,\cdot)$ 在 $\mathcal P_m^{\mathrm{loc}}\times\mathcal W^{\mathrm{loc}}$ 上联合连续。特别地，对每个固定的 $w\in\mathcal W^{\mathrm{loc}}$，映射 $p_m\mapsto S_m(p_m,w)$ 在 $\mathcal P_m^{\mathrm{loc}}$ 上连续。
\end{proposition}

\begin{proof}
先看需求对象。定义
\[
g_D(z,p_m):=x_B^*(z;p_m,M)\lambda_B(z,p_m,M).
\]
由 Assumption~\\ref{ass:p6:cm-cont}，$g_D$ 在紧集 $Z\times\mathcal P_m^{\mathrm{loc}}$ 上联合连续。根据 Extreme Value Theorem，连续函数在紧集上有界，所以存在有限常数 $K_D$，使得
\[
|g_D(z,p_m)|\le K_D
\quad\text{对所有 } (z,p_m)\in Z\times\mathcal P_m^{\mathrm{loc}}.
\]

现在取任意序列 $\{p_m^n\}_{n\ge 1}\subset\mathcal P_m^{\mathrm{loc}}$，满足 $p_m^n\to p_m$，其中 $p_m\in\mathcal P_m^{\mathrm{loc}}$。固定任意状态 $z\in Z$。由 $g_D$ 的联合连续性，
\[
g_D(z,p_m^n)\to g_D(z,p_m)
\quad\text{当 } n\to\infty.
\]
同时，对每个 $n$ 与每个 $z\in Z$ 都有
\[
|g_D(z,p_m^n)|\le K_D.
\]
因为 $\mu_B$ 是有限测度，常数函数 $z\mapsto K_D$ 关于 $\mu_B$ 可积：
\[
\int_Z K_D\,\mu_B(dz)=K_D\mu_B(Z)<\infty.
\]
于是可以应用 Dominated Convergence Theorem，得到
\[
\lim_{n\to\infty} D_m(p_m^n,M)
=
\lim_{n\to\infty}\int_Z g_D(z,p_m^n)\,\mu_B(dz)
=
\int_Z g_D(z,p_m)\,\mu_B(dz)
=
D_m(p_m,M).
\]
这就证明了 $D_m(\cdot,M)$ 在 $\mathcal P_m^{\mathrm{loc}}$ 上连续。

再看供给对象。定义
\[
g_S(z,p_m,w):=m^*(z;p_m,w).
\]
由 Assumption~\\ref{ass:p6:cm-cont}，$g_S$ 在紧集 $Z\times\mathcal P_m^{\mathrm{loc}}\times\mathcal W^{\mathrm{loc}}$ 上联合连续。再次由 Extreme Value Theorem，存在有限常数 $K_S$，使得
\[
|g_S(z,p_m,w)|\le K_S
\quad\text{对所有 } (z,p_m,w)\in Z\times\mathcal P_m^{\mathrm{loc}}\times\mathcal W^{\mathrm{loc}}.
\]

取任意序列 $\{(p_m^n,w^n)\}_{n\ge 1}\subset\mathcal P_m^{\mathrm{loc}}\times\mathcal W^{\mathrm{loc}}$，满足
\[
(p_m^n,w^n)\to(p_m,w),
\]
其中 $(p_m,w)\in\mathcal P_m^{\mathrm{loc}}\times\mathcal W^{\mathrm{loc}}$。固定任意状态 $z\in Z$。由 $g_S$ 的联合连续性，
\[
g_S(z,p_m^n,w^n)\to g_S(z,p_m,w)
\quad\text{当 } n\to\infty.
\]
同时，对每个 $n$ 与每个 $z\in Z$，
\[
|g_S(z,p_m^n,w^n)|\le K_S.
\]
由于 $\mu_C$ 是有限测度，
\[
\int_Z K_S\,\mu_C(dz)=K_S\mu_C(Z)<\infty.
\]
因此可以应用 Dominated Convergence Theorem，得到
\[
\lim_{n\to\infty} S_m(p_m^n,w^n)
=
\lim_{n\to\infty}\int_Z g_S(z,p_m^n,w^n)\,\mu_C(dz)
=
\int_Z g_S(z,p_m,w)\,\mu_C(dz)
=
S_m(p_m,w).
\]
于是 $S_m(\cdot,\cdot)$ 在 $\mathcal P_m^{\mathrm{loc}}\times\mathcal W^{\mathrm{loc}}$ 上联合连续。对固定 $w$ 时的 $p_m$ 连续性，就是这个联合连续性的特例。
\end{proof}

\paragraph{Market-clearing restriction and excess demand.}
令 $\mu:=(\mu_B,\mu_C)$ 表示进入 contract-manufacturing market 的 candidate measure pair。benchmark 的 market-clearing restriction 可以写成
\[
D_m(p_m,M)-S_m(p_m,w)=0.
\]
相应地，定义 contract-manufacturing excess-demand object 为
\[
Z_m(p_m,w,M;\mu):=D_m(p_m,M)-S_m(p_m,w).
\]
在 Assumption~\\ref{ass:p6:cm-meas} 与 Assumption~\\ref{ass:p6:cm-bound} 下，这个对象对固定 candidate prices 与固定 candidate measures 良定义。在 Assumption~\\ref{ass:p6:cm-cont} 下，它在紧局部价格域上连续。本部分不求解方程 $Z_m(p_m,w,M;\mu)=0$，也不声称 clearing price 已经存在或唯一。在当前 Phase 6 阶段，这里只构造 benchmark 的 market-clearing restriction 与 associated excess-demand object，供后续 equilibrium analysis 使用。

下一部分将按同样的 conditional logic 推导 labor-market object。

\subsection{Part 2：劳动力市场}

\paragraph{Benchmark object restatement.}
benchmark labor-market block 由 aggregate labor-demand object
\[
\begin{aligned}
L_D(M)={}&\int_Z \ell_A(z;w)\,\mu_A(dz)
+\int_Z \ell_B(z;p_m,M)\,\mu_B(dz)\\
&+\int_Z \ell_C(z;p_m,w)\,\mu_C(dz)
+\sum_{j\in\{A,B,C\}}\ell_E^j n_E^j(M),
\end{aligned}
\]
以及 labor-market-clearing restriction
\[
L_D(M)=\bar L
\]
来刻画。这里保持 benchmark object 的原始写法，不把它替换成新的公开记号。在当前 Phase 6 中，这个对象是在 candidate prices $(p_m,w)$、candidate finite Borel measures $(\mu_A,\mu_B,\mu_C)$ 和来自 Phase 5 的 entrant masses 下被评价的，但正文仍保留 benchmark 的记号 $L_D(M)$。

\paragraph{Economic role.}
$L_D(M)$ 中的每一项都有直接的经济含义。第一项是 incumbent type-$A$ labor demand，对应 integrated firms 的劳动需求并对其 candidate measure 积分。第二项是 incumbent type-$B$ labor demand，对应 research-specialized firms 的劳动需求并对其 candidate measure 积分。第三项是 incumbent type-$C$ labor demand，对应 manufacturing-specialized firms 的劳动需求并对其 candidate measure 积分。最后一项是 entrant setup labor：每一类 entrant $j\in\{A,B,C\}$ 都需要 $\ell_E^j$ 单位的 setup labor，因此 entrant side 的总劳动投入就是 $\ell_E^j n_E^j(M)$ 在三类 entrant 上的加总。

这里最重要的 fidelity point 是：$\ell_A(z;w)$、$\ell_B(z;p_m,M)$ 和 $\ell_C(z;p_m,w)$ 不是在本节新建出来的 factor-demand problem 的解。它们应当被理解为从已经优化完成的 flow-payoff blocks 中恢复出来的 labor-demand objects。本节所做的，只是把这些已存在的 labor objects 结合 candidate incumbent measures 与 Phase 5 的 entrant masses 聚合起来，从而得到 industry-level labor demand。

\paragraph{Mathematical problem.}
固定制度 $M\in\{0,1\}$、candidate prices $(p_m,w)$、candidate finite Borel measures $(\mu_A,\mu_B,\mu_C)$，以及 entrant masses $\bigl(n_E^A(M),n_E^B(M),n_E^C(M)\bigr)$。本部分要完成的工作是克制的。我们只需要说明：在这些 candidate inputs 下，benchmark labor-demand object $L_D(M)$ 是否良定义且有限；在局部紧价格域上，如果 labor objects 对价格连续，那么 aggregate labor demand 是否也随价格连续变化；最后，再把 labor-market-clearing restriction 写成一个 labor excess-demand object，供后续 equilibrium analysis 使用。本部分不求解 clearing wage，也不声称 clearing wage 唯一。

\paragraph{最小正则条件。}
为了使 benchmark labor object 可以用于后续分析，我们只需要很少的补充结构。第一，对固定的 candidate prices，三个 labor-demand objects $\ell_A(\cdot;w)$、$\ell_B(\cdot;p_m,M)$ 与 $\ell_C(\cdot;p_m,w)$ 必须关于状态 $z$ Borel 可测，因为它们要分别对相应的 candidate measure 积分。第二，对固定的 candidate prices，这三个 labor-demand objects 必须在紧状态空间 $Z$ 上有界，这样对有限测度积分时才一定得到有限值。第三，entrant coefficients $\ell_E^j$ 与 entrant masses $n_E^j(M)$ 都必须有限，这样 entrant labor term 才是有限和。最后，如果我们希望 aggregate labor demand 在局部紧价格域上连续，就还需要三个 labor-demand objects 对其相关价格参数局部连续。

为了后文引用方便，把这些条件记为：
\begin{assumption}[可测性与有限对象]\label{ass:p6:lm-meas}
对每个固定的三元组 $(p_m,w,M)$，映射 $z\mapsto \ell_A(z;w)$、$z\mapsto \ell_B(z;p_m,M)$ 与 $z\mapsto \ell_C(z;p_m,w)$ 在 $Z$ 上 Borel 可测。候选测度 $\mu_A,\mu_B,\mu_C$ 是定义在 $Z$ 上的有限 Borel measures。entrant coefficients $\ell_E^A,\ell_E^B,\ell_E^C$ 与 entrant masses $n_E^A(M),n_E^B(M),n_E^C(M)$ 都是有限的。
\end{assumption}

\begin{assumption}[固定候选价格下的劳动需求有界性]\label{ass:p6:lm-bound}
对每个固定的三元组 $(p_m,w,M)$，存在有限常数 $B_A(w)$、$B_B(p_m,M)$ 与 $B_C(p_m,w)$，使得
\[
|\ell_A(z;w)|\le B_A(w),\qquad
|\ell_B(z;p_m,M)|\le B_B(p_m,M),\qquad
|\ell_C(z;p_m,w)|\le B_C(p_m,w)
\]
对所有 $z\in Z$ 成立。
\end{assumption}

\begin{assumption}[候选价格的局部连续性]\label{ass:p6:lm-cont}
固定制度 $M$。令 $\mathcal P_m^{\mathrm{loc}}\subset\mathbb R_+$ 与 $\mathcal W^{\mathrm{loc}}\subset\mathbb R_+$ 表示紧的局部价格域。映射 $(z,w)\mapsto \ell_A(z;w)$ 在 $Z\times\mathcal W^{\mathrm{loc}}$ 上联合连续，映射 $(z,p_m)\mapsto \ell_B(z;p_m,M)$ 在 $Z\times\mathcal P_m^{\mathrm{loc}}$ 上联合连续，映射 $(z,p_m,w)\mapsto \ell_C(z;p_m,w)$ 在 $Z\times\mathcal P_m^{\mathrm{loc}}\times\mathcal W^{\mathrm{loc}}$ 上联合连续。
\end{assumption}

\begin{proposition}[总劳动需求的良定义性]
固定制度 $M$、candidate prices $(p_m,w)$、candidate finite Borel measures $(\mu_A,\mu_B,\mu_C)$ 以及 entrant masses $\bigl(n_E^A(M),n_E^B(M),n_E^C(M)\bigr)$。在 Assumption~\\ref{ass:p6:lm-meas} 与 Assumption~\\ref{ass:p6:lm-bound} 下，benchmark aggregate labor-demand object $L_D(M)$ 良定义且有限。
\end{proposition}

\begin{proof}
把 $L_D(M)$ 分成四部分：
\[
L_D(M)=L_A+L_B+L_C+L_E,
\]
其中
\[
L_A:=\int_Z \ell_A(z;w)\,\mu_A(dz),\quad
L_B:=\int_Z \ell_B(z;p_m,M)\,\mu_B(dz),\quad
L_C:=\int_Z \ell_C(z;p_m,w)\,\mu_C(dz),
\]
以及
\[
L_E:=\sum_{j\in\{A,B,C\}}\ell_E^j n_E^j(M).
\]
由 Assumption~\\ref{ass:p6:lm-meas}，三个 incumbent labor-demand objects 关于 $z$ 都是 Borel 可测，因此三个积分都是合法的可测积分。由 Assumption~\\ref{ass:p6:lm-bound}，
\[
|\ell_A(z;w)|\le B_A(w)\quad\text{对所有 } z\in Z.
\]
于是
\[
\int_Z |\ell_A(z;w)|\,\mu_A(dz)\le B_A(w)\mu_A(Z)<\infty,
\]
因为 $\mu_A$ 由 Assumption~\\ref{ass:p6:lm-meas} 是有限测度。完全平行地，
\[
\int_Z |\ell_B(z;p_m,M)|\,\mu_B(dz)\le B_B(p_m,M)\mu_B(Z)<\infty
\]
以及
\[
\int_Z |\ell_C(z;p_m,w)|\,\mu_C(dz)\le B_C(p_m,w)\mu_C(Z)<\infty.
\]
因此三个 incumbent components 都是有限的。

对 entrant term，Assumption~\\ref{ass:p6:lm-meas} 说明每个 coefficient $\ell_E^j$ 和每个 entrant mass $n_E^j(M)$ 都有限，所以每个乘积 $\ell_E^j n_E^j(M)$ 都有限。由于 entrant types 只有三类，$L_E$ 是有限个有限数的和，因此也是有限的。四部分合起来就得到 $L_D(M)$ 良定义且有限。
\end{proof}

\begin{proposition}[紧局部价格域上的连续性]
固定制度 $M$、candidate finite Borel measures $(\mu_A,\mu_B,\mu_C)$ 以及 entrant masses $\bigl(n_E^A(M),n_E^B(M),n_E^C(M)\bigr)$。在 Assumption~\\ref{ass:p6:lm-meas}--Assumption~\\ref{ass:p6:lm-cont} 下，benchmark labor-demand object $L_D(M)$ 关于 $(p_m,w)$ 在 $\mathcal P_m^{\mathrm{loc}}\times\mathcal W^{\mathrm{loc}}$ 上联合连续。
\end{proposition}

\begin{proof}
entrant term 在本阶段不随 candidate prices 变化，所以连续性问题只来自三个 incumbent integrals。由 Assumption~\\ref{ass:p6:lm-cont}，三个 labor-demand objects 在各自相关的紧域上联合连续。根据 Extreme Value Theorem，它们在这些紧域上都一致有界。由于候选测度都是有限测度，这些一致上界就是可积的 dominating functions。

现在取任意收敛序列 $(p_m^n,w^n)\to(p_m,w)$。由 Assumption~\\ref{ass:p6:lm-cont} 的联合连续性，三个 labor-demand objects 在每个固定状态 $z$ 上都逐点收敛到对应极限值。再由上面的统一有界性和候选测度有限性，可以分别对三个积分应用 Dominated Convergence Theorem。于是三个 incumbent integrals 都连续，entrant term 又是常数项，所以 $L_D(M)$ 关于 $(p_m,w)$ 联合连续。
\end{proof}

\paragraph{劳动市场出清限制与劳动超额需求。}
benchmark 的 labor-market-clearing restriction 是
\[
L_D(M)=\bar L.
\]
为了后续 equilibrium analysis，把同一个条件记成 labor excess demand 更方便：
\[
Z_\ell(p_m,w,M;\mu,n_E):=L_D(M)-\bar L,
\]
其中 $\mu:=(\mu_A,\mu_B,\mu_C)$，而 $n_E:=\bigl(n_E^A(M),n_E^B(M),n_E^C(M)\bigr)$。在 Assumption~\\ref{ass:p6:lm-meas} 与 Assumption~\\ref{ass:p6:lm-bound} 下，这个对象对 candidate prices、candidate measures 与 entrant masses 良定义；在 Assumption~\\ref{ass:p6:lm-cont} 下，它在紧局部价格域上连续。但在当前阶段，它仍然只是供后续 equilibrium analysis 使用的 market restriction 与 excess-demand object，而不是 clearing wage 的存在性或唯一性结论。

Phase 6 剩下的部分将把这些 conditional market objects 总结成后续 equilibrium analysis 可以直接调用的接口。

\subsection{Part 3：条件市场限制与阶段交接}

\paragraph{汇总市场对象。}
令
\[
\mu:=(\mu_A,\mu_B,\mu_C)
\qquad\text{以及}\qquad
n_E:=\bigl(n_E^A(M),n_E^B(M),n_E^C(M)\bigr).
\]
在给定 candidate prices $(p_m,w)$、candidate incumbent measures $\mu$ 和 entrant masses $n_E$ 的条件下，Phase 6 已经构造出 contract-manufacturing excess-demand object
\[
Z_m(p_m,w,M;\mu):=D_m(p_m,M)-S_m(p_m,w)
\]
以及 labor excess-demand object
\[
Z_\ell(p_m,w,M;\mu,n_E):=L_D(M)-\bar L.
\]
相应的 market-clearing restrictions 就是
\[
Z_m(p_m,w,M;\mu)=0
\qquad\text{以及}\qquad
Z_\ell(p_m,w,M;\mu,n_E)=0.
\]
这些都只是条件性的市场限制。在当前阶段，Phase 6 还没有证明同一组 candidate incumbent measures 会在 optimal behavior、exit 和 entrant inflows 之下满足 invariant-measure 条件。

\paragraph{Phase 6 已经建立了什么。}
Phase 6 已经建立了 benchmark contract-manufacturing market objects 的定义，建立了 benchmark aggregate labor-demand object 的定义，也把两个 market-clearing restrictions 都写成了 excess-demand notation。具体地说，$Z_m(p_m,w,M;\mu)$ 与 $Z_\ell(p_m,w,M;\mu,n_E)$ 现在都已经可作为模型市场侧的紧凑表示，供后续 equilibrium analysis 使用。同样重要的是，Phase 6 还没有建立什么：它还没有给出 invariant-measure result，还没有给出 stationary competitive equilibrium existence result，也还没有求出 values、policy rules、candidate measures、entrant masses 与 prices 上的联合 fixed point。

\paragraph{对后续阶段的交接。}
Phase 7 必须把 selected policy behavior、transition kernels、exit 与 entrant inflows 组合起来，写出 incumbent measures 的 law of motion；只有完成这一步，invariant measures 才能被定义和研究。Phase 8 则必须把 value functions、operating controls、organizational choices、prices、entrant masses 与 invariant measures 组装成 stationary competitive equilibrium 的定义及其 existence framework。
